% from ega1-7-fr.tex
\oldpage[I]{156}
à laquelle appartient $f$, on définit une application canonique de
$\operatorname{Hom}_S(X,Y)$ dans l'ensemble des $S$-applications rationnelles
de $X$ dans $Y$. On désigne par $\Gamma_{\mathrm{rat}}(X/Y)$ l'ensemble des
$Y$-sections rationnelles de $X$, et on a donc une application canonique
$\Gamma(X/Y)\to\Gamma_{\mathrm{rat}}(X/Y)$. Il est clair en outre que si $X$ et
$Y$ sont deux $S$-préschémas, l'ensemble des $S$-applications rationnelles de
$X$ dans $Y$ s'identifie canoniquement à
$\Gamma_{\mathrm{rat}}((X\times_S Y)/X)$
(\hyperref[I.3.3.14-fr]{3.3.14}).
\end{definition}

\begin{env}[7.1.3]
\label{I.7.1.3-fr}
Il résulte aussitôt de (\hyperref[I.7.1.2-fr]{7.1.2}) et de
(\hyperref[I.3.3.14-fr]{3.3.14}) que les fonctions rationnelles sur $X$
s'identifient canoniquement aux
\emph{classes d'équivalence des sections du faisceau structural
$\mathscr{O}_X$} au-dessus d'ouverts partout denses de $X$, deux telles
sections étant équivalentes si elles coïncident dans un ouvert partout dense
contenu dans l'intersection de leurs ensembles de définition. Il en résulte en
particulier que les fonctions rationnelles sur $X$ forment un
\emph{anneau} $R(X)$.
\end{env}

\begin{env}[7.1.4]
\label{I.7.1.4-fr}
Lorsque $X$ est un préschéma \emph{irréductible}, tout ouvert non vide est dense
dans $X$~; on peut encore dire que les ouverts non vides de $X$ sont les
\emph{voisinages ouverts du point générique $x$} de $X$. Dire que deux
morphismes de parties ouvertes non vides de $X$ dans $Y$ sont équivalents
signifie donc dans ce cas qu'ils ont \emph{même germe} au point $x$. Autrement
dit, les applications rationnelles (resp. $S$-applications rationnelles)
$X\to Y$ s'identifient aux \emph{germes de morphismes} (resp. de
$S$-morphismes) de parties ouvertes non vides de $X$ dans $Y$ au point
générique $x$ de $X$. En particulier~:
\end{env}

\begin{proposition}[7.1.5]
\label{I.7.1.5-fr}
Si $X$ est un préschéma irréductible, l'anneau $R(X)$ des fonctions rationnelles
sur $X$ s'identifie canoniquement à l'anneau local $\mathscr{O}_x$ du point
générique $x$ de $X$. C'est un anneau local de dimension 0, et par suite un
anneau local artinien lorsque $X$ est noethérien~; c'est un corps lorsque $X$
est intègre, et il s'identifie au corps des fractions de $A(X)$ lorsqu'en outre
$X$ est un schéma affine.
\end{proposition}

Vu ce qui précède et l'identification des fonctions rationnelles aux sections
de $\mathscr{O}_X$ au-dessus d'un ouvert partout dense, la première assertion
n'est autre que la définition de la fibre d'un faisceau en un point. Pour les
autres assertions, on peut se borner au cas où $X$ est affine d'anneau $A$~;
alors $\mathfrak{j}_x$ est le nilradical de $A$, et $\mathscr{O}_x$ est donc de
dimension 0~; si $A$ est intègre, $\mathfrak{j}_x=(0)$, et $\mathscr{O}_x$ est
donc le corps des fractions de $A$. Enfin, si $A$ est noethérien, on sait
(\cite{I-11}, p.~127, cor.~4) que $\mathfrak{j}_x$ est nilpotent et
$\mathscr{O}_x=A_{\mathfrak{j}_x}$ artinien.

Si $X$ est \emph{intègre}, l'anneau $\mathscr{O}_z$ est intègre pour tout
$z\in X$~; tout ouvert affine $U$ contenant $z$ contient aussi $x$, et $R(U)$,
égal au corps des fractions de $A(U)$, s'identifie donc à $R(X)$~; on en
conclut que $R(X)$ s'identifie aussi au
\emph{corps des fractions de $\mathscr{O}_z$}~: l'identification canonique de
$\mathscr{O}_z$ à un sous-anneau de $R(X)$ consiste à faire correspondre à tout
germe de section $s\in\mathscr{O}_z$ l'unique fonction rationnelle sur $X$,
classe d'une section de $\mathscr{O}_X$ (nécessairement définie sur un ouvert
partout dense) ayant pour germe $s$ au point $z$.

\begin{env}[7.1.6]
\label{I.7.1.6-fr}
Supposons maintenant que $X$ ait un nombre \emph{fini} de composantes
irréductibles $X_i$ ($1\leq i\leq n$) (ce qui est le cas lorsque l'espace
sous-jacent à $X$ est \emph{noethérien})~; soit $X'_i$ l'ouvert de $X$,
complémentaire par rapport à $X_i$ de la réunion des $X_j\cap X_i$ pour
$j\neq i$~; $X'_i$ est irréductible, son point générique $x_i$ est le point
générique de $X_i$, et les $X'_i$ sont deux à deux disjoints, leur réunion
étant dense dans $X$ (\hyperref[0.2.1.6-fr]{0, 2.1.6}). Pour
